\section{Conservation and control of the defect}\label{revision:conservation}
The equality case in Section~\ref{ms:geometry-counts} identifies the
extremal Bernoulli law. To connect that geometry to the excess, we
integrate the nonnegative contact difference $\Delta=Q-h$ against the
zero-bias density. This integral has two useful descriptions. Since
$\Delta$ vanishes on the support, integration over the gaps produces
$J$ and a nonnegative leakage term. The identity
$\int\tau g'=\E[Xg(X)]$ gives a second description in terms of
moments and the sum distribution. Equating them yields conservation.
The coefficient $\CE=(3+\sqrt{10})\phi(0)/6$ then combines the
third-moment penalty with the gap term into the nonnegative defect $D$.

Retain the definitions of $a,Q,h$ and $\Delta=Q-h$ from
Section~\ref{v9:core:section}. For now let $(P,H)$ be an attained pair
at any $c>0$, so the law has compact support $K$ and $\Delta=0$ on $K$.

\subsection{Zero-bias integration and the conservation identity}

Put $\nu=\mathcal L(H-S_{n-1})$ and
$\tau(t)=\E[X\mathbf1_{X\ge t}]$. Then $\tau\ge0$, $\int\tau=1$.
Over the bounded complementary intervals $(u,v)$ of the closed support
$K$, define
\begin{align*}
 J&=\sum\tau_{u,v}(v-u)^2,&
 J_1&=\sum\tau_{u,v}\int_u^v(t-u)t\,dt,\\
 J_2&=\sum\tau_{u,v}\int_u^v(t-u)t|t|\,dt,&
 L&=\sum\tau_{u,v}\int_{[u,v)}(t-u)\,d\nu(t).
\end{align*}
Here $\tau_{u,v}$ is the constant interior value; $J$ and $\tau$ agree
with the definitions in Section~\ref{ms:geometry-counts}. The leakage
$L$ is a weighted contribution of the law of $H-S_{n-1}$ inside support
gaps, and is nonnegative because its measure and integrand are
nonnegative. This sign needs no estimate for $C_{n-1}$.
The integral uses $[u,v)$: the right endpoint belongs to
$h(v)$, whereas an atom at the left endpoint has zero weight. If
$K\subset[-R,R]$, then $J\le2R$, $|J_1|\le RJ/2$,
$|J_2|\le R^2J/2$, and $0\le L\le2R$, so all gap sums converge.

The remainder has four parts: a constant moment and gap term, a linear
threshold term, a quadratic threshold term, and the Fourier error $E_n$.
Specifically, let $\phi_z=\phi(H/\sqrt n)$,
$M_4=\E X^4$, $M_5=\E|X|^5$, and put
\[
 K_q=M_5/4-\rho M_4/2-3\beta\gamma/2+\rho/2,
 \qquad K_1=\E[X|X|^3]-3\beta-3\rho\gamma/2,
\]
\[
 \Omega=J_1+(M_4-3)/6,\quad
 \varepsilon=\E q_z(S_n/\sqrt n)-\E q_z(Z_0)
                  -2\gamma\phi_z/(3\sqrt n),\quad
 q_z(x)=(x-z)|x-z|,
\]
where $z=H/\sqrt n$ and $Z_0$ is standard normal. Define $E_n=-n^{3/2}\varepsilon/2$ and
\begin{equation}\label{v9:core:R}
 \mathcal R_c=c(3\beta J/2+3\rho J_1-3J_2-K_q)
 +H(cK_1+\phi_z\Omega)-H^2(c\rho+\gamma\phi_z/2)+E_n.
\end{equation}


\begin{theorem}[Conservation at contact]\label{revision:contact-conservation}
At an attained pair with $c=\CE$, put $e=\exp(-H^2/(2n))$. Then
\begin{equation}\label{v9:core:positive}
 \mathcal R_c=(n+H^2)a+nc\rho(1-e)+\frac{n\phi(0)e}{6}D+n^{3/2}L.
\end{equation}
More generally, at an attained pair for any $c>0$,
\begin{equation}\label{v9:core:conservation}
 (n+H^2)a+nc\rho-\frac{n\phi_z}{6}(\gamma+3J)+n^{3/2}L=\mathcal R_c.
\end{equation}
\end{theorem}
The terms in \eqref{v9:core:positive} have separate roles:
\begin{center}\small
\begin{tabular}{@{}p{.40\linewidth}p{.54\linewidth}@{}}
\toprule
Term & Role \\
\midrule
$(n+H^2)a$ & The excess, weighted by $n+H^2$. \\
$nc\rho(1-e)$ & The loss from moving the normal threshold away from zero. \\
$n\phi(0)eD/6$ & The moment and gap defect from Section~\ref{ms:geometry-counts}. \\
$n^{3/2}L$ & A weighted contribution inside the support gaps. \\
$\mathcal R_c$ & The moment, threshold and Fourier terms to be bounded. \\
\bottomrule
\end{tabular}
\end{center}
The four terms on the right of \eqref{v9:core:positive} are nonnegative
when $a\ge0$. Solving for
the excess makes the role of the remainder explicit:
\begin{equation}\label{opus:excess-identity}
 a=\frac{\mathcal R_c-n\CE\rho(1-e)-n\phi(0)eD/6-n^{3/2}L}
             {n+H^2}.
\end{equation}
In particular, fix $T\ge0$, $H_0\ge0$ and $N\ge2$. For any such pair
with $a\ge0$, $n\ge N$, $\mathcal R_c\le T$ and $|H|\le H_0$, one has
\[
 D\le\frac{6T\exp(H_0^2/(2N))}{n\phi(0)}.
\]
Theorem~\ref{v9:entry:theorem} supplies such numerical bounds under
its stated hypotheses. Thus the $1/n$ control of the defect comes from
uniform quantitative estimates, rather than from a fixed-law limit.

\begin{proof}
\emph{1. Integration over the gaps.}
Let $\mathcal A=\int\tau(Q-h)$. Since $h(t)=\nu([t,\infty))$, one has
$h(t)=h(v)+\nu([t,v))$ inside a gap. Thus bounded Fubini gives
\[
 \int_u^v\Delta(t)dt=(v-u)\Delta(v)
 +\int_u^v(t-u)(-Q'(t))dt
 -\int_{[u,v)}(t-u)d\nu(t).
\]
An atom at $v$ belongs to $h(v)$ and is not in the last integral;
an atom at $u$ is included but multiplied by zero. Because $\Delta=0$ on the entire closed support, the endpoint terms
and the integral over $K$ vanish. Summing gives
\begin{equation}\label{v9:core:first}
 \mathcal A=\frac{\phi_z J}{2\sqrt n}
 +\frac{3c\beta J/2+(H\phi_z+3c\rho)J_1-3cJ_2}{n^{3/2}}-L.
\end{equation}

\emph{2. Zero-bias integration.} Signwise Fubini proves
$\int\tau g'=\E[Xg(X)]$ for absolutely continuous $g$ with bounded
derivative on the relevant compact range. In particular
\[
 \int t\tau=\gamma/2,\quad \int t^2\tau=M_4/3,
 \quad\int |t|^3\tau=M_5/4.
\]
Apply that same identity conditionally to
$g_s(x)=\min(s+x,H)$ and then use exchangeability. The choice of
derivative at its single corner is immaterial under the density $\tau$.
It follows directly that
\[
 \int\tau(t)h(t)dt=\E[X_n\min(S_n,H)]
 =\frac1n\E[S_n\min(S_n,H)]
 =T,
\]
where $T=\E[(Z^2-zZ)\mathbf1_{Z\le z}]$, $Z=S_n/\sqrt n$.
The last expression vanishes at an endpoint atom. Integrating
\eqref{v9:core:Q} against $\tau$ therefore yields
\begin{equation}\label{v9:core:second}
 \mathcal A=F-\frac{\gamma\phi_z}{2\sqrt n}
 +\frac{H\phi_z(3-M_4)/6+cK_q}{n^{3/2}}-T.
\end{equation}

\emph{3. The first-moment identity.} Set $M=\E[Z\mathbf1_{Z\le z}]$. Since
$\E[Xh(X)]=M/\sqrt n$, integration of $XQ(X)$ against $P$ gives
\begin{equation}\label{v9:core:firstmoment}
 M+\phi_z=-\frac{\gamma H\phi_z}{2n}+\frac{cK_1}{n}.
\end{equation}
\emph{4. Gaussian comparison and assembly.} The pointwise identity
$q_z(x)=(x-z)^2-2(x-z)^2\mathbf1_{x\le z}$, with the same identity
for a normal variable, gives
\[
 T-\Phi(z)=-\varepsilon/2-\gamma\phi_z/(3\sqrt n)
              +z(M+\phi_z)-z^2(F-\Phi(z)).
\]
Inserting \eqref{v9:core:firstmoment} and
$\sqrt n(F-\Phi(z))=c\rho+a$ into \eqref{v9:core:second} gives
\begin{align*}
 \sqrt n\mathcal A={}&c\rho+a-\gamma\phi_z/6+P_0/n
                    +\sqrt n\varepsilon/2,\\
 P_0={}&H^2(c\rho+a+\gamma\phi_z/2)
       -H\{cK_1+\phi_z(M_4-3)/6\}+cK_q.
\end{align*}
Equate this with \eqref{v9:core:first} multiplied by $\sqrt n$,
then multiply by $n$ and collect terms. This is
\eqref{v9:core:conservation}. Substituting
$c=\CE=(3+\sqrt{10})\phi(0)/6$, $\phi_z=\phi(0)e$, and
$\gamma+3J=(3+\sqrt{10})\rho-D$ proves
\eqref{v9:core:positive}.
\end{proof}

The compensated Fourier representation of this same $E_n$ is proved in
Appendix~\ref{v9:fourier:section}. Its comparator has transform
$e^{-nt^2/2}(1-in\gamma t^3/6)$. Appendix~\ref{revision:offcontact}
retains the stronger identity for arbitrary compactly supported laws,
including the residuals that disappear here by contact.

\subsection{Why the remainder bound must come first}
By Lemma~\ref{ms:physical-defect},
\begin{equation}\label{v9:core:defect}
 D=3(\rho+2m-J)+(\sqrt{10}\rho-\gamma-6m)\ge0.
\end{equation}
Specialize from now on to $c=\CE$ and $a\ge0$. The sharper estimate
below keeps the threshold loss. Its coefficient is not known positive
until a bound on $\mathcal R_c$ has first been established.

\begin{corollary}[From a remainder bound to a defect bound]\label{v9:core:license}\label{v9:core:phase-price}
At such a contact pair with $a\ge0$, assume $n\ge N\ge2$,
$\rho\ge l>0$, $0<c_-\le c$, and a proved remainder bound
$\mathcal R_c\le T<Nc_-l$. A negative $T$ excludes the pair.
Otherwise
\begin{gather*}
 \gamma+3J>0,\qquad e\ge1-T/(Nc_-l)>0,\\
 \mathcal U:=(n+H^2)a+n\phi(0)D/6+n^{3/2}L
       \le\mathcal R_c,\\
 \mathcal U\le\mathcal R_c-\tfrac12(c_-l-T/N)H^2.
\end{gather*}
No condition $D<2$ or own-score upper bound is required.
\end{corollary}
\begin{proof}
Set $M_c=c\rho-\phi(0)D/6=\phi(0)(\gamma+3J)/6$.
The contact identity gives
\[
 eM_c=(1+H^2/n)a+c\rho+\sqrt n L-\mathcal R_c/n
        \ge c\rho-T/n>0.
\]
Thus $M_c>0$ and, since $D\ge0$, $M_c\le c\rho$.
The density bound follows. Regrouping \eqref{v9:core:positive} gives
$\mathcal R_c=\mathcal U+n(1-e)M_c$, proving the unweighted bound.
Finally, put $w=H^2/(2n)$, so $e=e^{-w}$. Then
\[
 n(1-e)M_c=n(e^w-1)eM_c
 \ge\frac{H^2}{2}(c\rho-T/n)
 \ge\frac{H^2}{2}(c_-l-T/N).
\]
This is a bound on $\mathcal U$, not an
improved bound on the raw remainder.
\end{proof}

The corollary is used in two steps. First, an upper bound for
$\mathcal R_c$ proves $M_c>0$ and a positive lower bound for $e$.
Second, the now nonnegative term $n(1-e)M_c$ can be subtracted to
improve the bound for $\mathcal U$ by a multiple of $H^2$. This is the
improvement that helps bound the threshold. It is a bound on
$\mathcal U$, so it cannot be substituted for the initial upper bound
on $\mathcal R_c$. Every inequality holds separately at each $n,H$;
common bounds on a closed moment cell therefore also apply when the
maximizing law varies with $n$.
